\section{Retraction stability and signal transfer}
\label{sec:retraction_stability}

The signal floor (\Cref{lem:signal_floor}) is established on the
\emph{un-retracted} iterate $\tilde x_T$, where the running-average
identity $\sum_k w_{T,k}=1$ is available. The decoder, however, sees the
on-sphere iterate generated by the dynamics. Two on-sphere maps are in
play: the \emph{tangential-step retraction}
$x_t=\Retr_{x_{t-1}}(\Tproj{x_{t-1}}(\tilde x_t-x_{t-1}))$ of
\Cref{def:dynamics}, and the \emph{full normalization}
$\bar x_t\coloneqq\sqrt d\,\tilde x_t/\norm{\tilde x_t}_2$ that
\Cref{lem:running_average} attaches to the convex combination. These agree
only to second order; this lemma quantifies the gap. It is genuine work:
the normalization map is a second-order (not first-order, not isometric)
retraction, and the curvature correction must be quantified.

\begin{lemma}[Retraction stability and signal transfer]
\label{lem:retraction_stability}
Under \Cref{ass:bounded_architecture}, let $x_{t-1}\in\Sphere$, let
$\tilde x_t=(1-\alpha)x_{t-1}+\alpha V_t$ be the un-retracted iterate of
\Cref{def:dynamics}, let $x_t$ be its tangential-step retraction and
$\bar x_t=\sqrt d\,\tilde x_t/\norm{\tilde x_t}_2$ its full normalization.
Write the tangential step in unit-sphere coordinates as $s$ with
$\eta\coloneqq\norm{s}_2/\sqrt d$. Then $\Retr$ is a second-order
retraction (\cite{boumal2023}, Proposition~5.44, eq.~(5.30)), and for every
fixed token row $a\in\Vocab$,
\begin{align}\label{eq:retraction_transfer}
   \bigl|\, \inner{W_U^a}{x_t} - \inner{W_U^a}{\bar x_t} \,\bigr|
   \;\le\;
   C_{\mathrm{ret}}\, R_U\, \eta^3\,\sqrt d,
\end{align}
where $C_{\mathrm{ret}}$ is the absolute third-derivative bound of the
normalization map on a neighbourhood of $\Sphere$. Under the
radius-$\sqrt d$ rescaling $\eta=O(1/\sqrt d)$ of
\Cref{rem:retraction_rescaling}, the right-hand side is $O(1/d)$ on the
\emph{raw logit} and hence $O(1/d^{1.5})$ on the \emph{angular margin}
(after dividing by the row-norm scale $\norm{W_U^a}_2\sqrt d\asymp\sqrt d$).
Consequently any high-probability bound on a fixed-row logit established for
$\tilde x_T$ (equivalently for $\bar x_T$) transfers to the on-sphere
iterate $x_T$ up to an additive $O(1/d)$ logit error, dominated by the
noise scale $\sigma_T\asymp1/\sqrt{T_{\max}d}$ of \Cref{lem:signal_floor}
whenever $T_{\max}=O(d)$ (raw-logit comparison; $T_{\max}=O(d^2)$ in
angular-margin units).
\end{lemma}

\begin{proof}
The argument expands both on-sphere maps to second order on the unit
sphere, where they coincide, and bounds the residual; the radius-$\sqrt d$
rescaling then converts the unit-sphere residual to the logit scale. We
\emph{never} apply $\sum_k w_{T,k}=1$ to $x_t$; the representation lives on
$\tilde x_t$ (\Cref{lem:running_average}), and this lemma is the bridge to
the on-sphere iterate.

\smallskip\noindent\textbf{Step 1 (both maps are second-order retractions).}
Write $x=\sqrt d\,u$ with $u$ on the \emph{unit} sphere $S^{d-1}$, and let
$s\in T_u S^{d-1}$ be the unit-sphere tangential step, $\norm{s}_2=\eta$.
The unit-sphere normalization retraction is $\bar R_u(s)=(u+s)/\norm{u+s}_2$;
by \cite{boumal2023}, Example~3.49 and Proposition~5.44, it is a
\emph{second-order} retraction, with the Taylor expansion
\begin{align}\label{eq:retr_taylor}
   \bar R_u(s) \;=\; u + s - \tfrac12\norm{s}_2^2\,u + r(s),
   \qquad \norm{r(s)}_2\le C_{\mathrm{ret}}\,\eta^3,
\end{align}
where the cubic residual $r(s)$ is controlled by the third derivative of
$v\mapsto v/\norm{v}_2$ on a neighbourhood of $u$ (Boumal eq.~(5.30); valid
for $\eta$ below the injectivity radius, automatic for large $d$). The
tangential-step retraction $\Retr$ of \Cref{def:dynamics}, restricted to
the unit sphere, is the \emph{same} normalization map applied to the
tangential displacement, hence has the \emph{same} second-order expansion
\Cref{eq:retr_taylor} with its own cubic residual of the same order; the
$-\tfrac12\norm{s}_2^2u$ curvature term (radial, along $u$) is common to
both. Subtracting the two expansions, the linear and quadratic terms cancel
and only the cubic residuals remain:
\begin{align*}
   \norm{\,\Retr_{x_{t-1}}(\cdot)/\sqrt d - \bar R_u(s)\,}_2
   \;\le\; 2\,C_{\mathrm{ret}}\,\eta^3,
\end{align*}
on the unit sphere (we absorb the factor $2$ into $C_{\mathrm{ret}}$).

\smallskip\noindent\textbf{Step 2 (project onto a row, restore the radius).}
Restoring the radius, $x_t=\sqrt d\,\Retr_{x_{t-1}}(\cdot)/\sqrt d$ and
$\bar x_t=\sqrt d\,\bar R_u(s)$, so taking $\inner{W_U^a}{\cdot}$ and using
Cauchy--Schwarz with $\norm{W_U^a}_2\le R_U$,
\begin{align*}
   \bigl|\inner{W_U^a}{x_t}-\inner{W_U^a}{\bar x_t}\bigr|
   \;\le\; \norm{W_U^a}_2\cdot\sqrt d\cdot\norm{\,\Retr(\cdot)/\sqrt d-\bar R_u(s)\,}_2
   \;\le\; C_{\mathrm{ret}}\,R_U\,\eta^3\,\sqrt d,
\end{align*}
which is \Cref{eq:retraction_transfer}.

\smallskip\noindent\textbf{Step 3 (rescaling and accumulation).} With
$\eta=O(1/\sqrt d)$ (\Cref{rem:retraction_rescaling}),
$\eta^3\sqrt d=O(d^{-3/2}\cdot d^{1/2})=O(1/d)$ on the raw logit, and
$O(1/d^{1.5})$ once divided by the row-norm scale
$\norm{W_U^a}_2\sqrt d\asymp\sqrt d$ implicit in the angular margin. The
\emph{operational} iterate read by the verifier is the attention average
followed by a \emph{single} LayerNorm, which is exactly
$\bar x_T=\sqrt d\,\tilde x_T/\norm{\tilde x_T}$ of
\Cref{eq:onsphere_direction}; for this iterate \emph{there is no accumulated
retraction error}, and the entire analysis (which lives on $\tilde x_T$)
applies without correction. The lemma's content is the comparison to the
\emph{alternative} per-step reading $x_t=\Retr_{x_{t-1}}(\cdots)$ of
\Cref{def:dynamics}: \Cref{eq:retraction_transfer} bounds the single-step
discrepancy by $O(1/d)$ (raw logit). \textbf{Domination.} Taking the
operational iterate $\bar x_T$, the only geometric error is this $O(1/d)$
single-step quantity, dominated by the noise scale
$\sigma_T\asymp R_U M e^S/\sqrt{T_{\max}d}$ of \Cref{lem:signal_floor}
whenever $1/d\le1/\sqrt{T_{\max}d}$, i.e.\ $T_{\max}\le d$; in
angular-margin units the error is $O(1/d^{1.5})$ and the condition relaxes
to $T_{\max}\le d^2$. (Should one instead track the per-step-retraction
trajectory of \Cref{def:dynamics}, the gaps could accumulate to $O(T_{\max}/d)$;
we avoid this by analysing the operationally-exact single-normalization
iterate $\bar x_T$.) This completes the proof.
\end{proof}

\begin{remark}[The radius-$\sqrt d$ rescaling: why $\eta=O(1/\sqrt d)$]
\label{rem:retraction_rescaling}
Writing $x=\sqrt d\,u$ with $u$ on the \emph{unit} sphere, a single
SGDM step displaces the unit direction by $s=\Tproj{x_{t-1}}(\tilde x_t-x_{t-1})/\sqrt d$,
the tangential part of $\tilde x_t-x_{t-1}=\alpha(V_t-x_{t-1})$ rescaled by
the radius. Because $\norm{\alpha(V_t-x_{t-1})}_2\le\alpha(M+\sqrt d)=O(\sqrt d)$
and the tangential projection only shrinks the norm, the induced
unit-sphere displacement is $\eta=\norm{s}_2/\sqrt d=O(1)$ in the worst
case; but the \emph{signal-carrying} part that moves the angle is
$O(1)$ in ambient units against a radius $\sqrt d$, so the per-step angular
increment is $\eta=O(1/\sqrt d)$. Cubing gives $\eta^3=O(d^{-3/2})$, the
source of the headline. The retraction residual is deterministic and
geometric; it is a \emph{separate} error source from the martingale noise
of \Cref{lem:azuma}, and the two add. We state the raw-logit exponent
$O(1/d)$ and the angular-margin exponent $O(1/d^{1.5})$ honestly, rather
than claiming the smaller one for the logit; either way the error is
dominated by the noise in the stated horizon regime.
\end{remark}
